\chapter*{摘\quad 要}
\addcontentsline{toc}{chapter}{摘要}

本项目面向工业巡检场景，设计并实现了一套全栈国产化的自主巡检机器人系统，用于代替人工完成区域遍历与工业仪表读数分析。系统以开源鸿蒙（OpenHarmony 5.0）为操作系统、以华为昇腾 310B NPU 为边缘算力，不使用 ROS 等境外机器人中间件，从硬件搭建、激光雷达驱动到 SLAM 建图、路径规划、仪表识别与大模型分析均自主实现，数据不出本体、断网可用。

系统采用三节点异构架构：每辆车由一块紫派和一块香橙派组成，紫派（OpenHarmony 5.0）负责激光感知、SLAM 导航与运动控制，香橙派（昇腾 310B）负责仪表检测、关键点定位与读数几何换算，并离线运行 DeepSeek 大模型生成巡检分析报告；香橙派识别到仪表后联动紫派减速，读数完成后恢复行进。鸿蒙平板（ArkTS）负责人机交互、实时可视化与多车总控。多机协同时，各车与平板通过华为分布式软总线共享同一份任务状态，按划分的子区域协同完成大区域全覆盖巡检。

实测表明，系统可在室内场景自主建图、导航避障与全覆盖遍历；仪表识别端到端延迟约 34~ms（约 29~FPS），支持一个画面内多块仪表同时读数与越限告警；端侧大模型可全程离线生成趋势分析报告，验证了国产软硬件承载完整机器人系统的可行性。

\vspace{1.5em}
\noindent\textbf{关键词：}开源鸿蒙（OpenHarmony）；分布式软总线；昇腾边缘计算；同步定位与建图（SLAM）；全覆盖路径规划；工业仪表识别
